% =====================================================================
%  Section: Conclusion
% =====================================================================
\section{Conclusion}
\label{sec:conclusion}

We presented \textbf{Canvas3D}, a four-module framework that transforms static
oil paintings and ink-wash artworks into 3D-aware, artistically controllable
animated videos. Our key design insight is the \emph{modular separation of
perceptual, geometric, compositional, and temporal concerns} across four
specialized modules, each of which can be optimized and upgraded independently
and which communicate only through small, interpretable handshake
representations.

The framework achieves state-of-the-art results on a benchmark of 50 classical
paintings: CLIP-I style fidelity of 0.86, LPIPS perceptual quality of 0.28,
disocclusion PSNR of 29.1\,dB, and FID-VID temporal coherence of 18.3. User
studies confirm strong preference for Canvas3D outputs across artistic
authenticity, motion naturalness, and compositional-intent alignment.
Ablations validate each module's contribution, with the largest degradations
occurring when the depth-aware inpainting module or the edge control branch
is removed.

\smallskip
\noindent\textbf{Limitations.}
Several limitations suggest directions for future work:
(i)~the per-frame conditioning operates at the 2D projection level, so
extreme viewpoint changes (>$60^\circ$ rotation) introduce hallucinated
textures inconsistent with the underlying 3D model;
(ii)~the feed-forward 3D predictor's single-image reconstruction quality
depends on the object's visibility in the original painting---heavily
occluded regions produce degraded 3D geometry;
(iii)~our color-inpainting module assumes locally similar textures,
which may fail for large disocclusion regions in sparse compositions;
(iv)~the current pipeline processes objects serially; parallelized
scheduling for multiple 3D lifting calls would reduce end-to-end latency.

\smallskip
\noindent\textbf{Future work.}
We plan to explore: (a)~multi-canvas scene fusion, enabling composition
from objects extracted from different paintings; (b)~neural radiance
field representations for the background layer, which would allow more
flexible camera trajectories; (c)~interactive real-time line-art
composition with progressive 3D rendering; and (d)~audio-visual
alignment for narratively-scored animation generation.
